\documentclass[11pt,a4paper]{article}
\usepackage[utf8]{inputenc}
\usepackage[T1]{fontenc}
\usepackage[spanish,es-nodecimaldot]{babel}
\usepackage{lmodern}
\usepackage[a4paper,margin=2.2cm]{geometry}
\usepackage{xcolor}
\usepackage{hyperref}
\usepackage{enumitem}
\usepackage{booktabs}
\usepackage{longtable}
\usepackage{tabularx}
\usepackage{array}
\usepackage{listings}
\usepackage[most]{tcolorbox}
\usepackage{fancyhdr}
\usepackage{microtype}
\usepackage[strings]{underscore}
\usepackage{textcomp}

\newcommand{\NombreEstudiante}{Nombre: \rule{6cm}{0.4pt}}
\newcommand{\FechaPrueba}{Por confirmar}
\newcommand{\MinutosSesion}{40}
\newif\ifmostrarSoluciones
\mostrarSolucionestrue

\definecolor{azulUCM}{HTML}{17365D}
\definecolor{azulClaro}{HTML}{EAF2F8}
\definecolor{verdeClaro}{HTML}{E9F7EF}
\definecolor{amarilloClaro}{HTML}{FFF8DC}
\definecolor{rojoClaro}{HTML}{FDEDEC}
\definecolor{grisCodigo}{HTML}{F4F6F7}
\definecolor{sqlKeyword}{HTML}{005CC5}
\definecolor{sqlString}{HTML}{A31515}
\definecolor{sqlComment}{HTML}{22863A}

\hypersetup{colorlinks=true,linkcolor=azulUCM,urlcolor=blue,
  pdftitle={Guía de lectura activa -- Unidad 2 de Bases de Datos},
  pdfauthor={Material de estudio INF-324}}
\pagestyle{fancy}\fancyhf{}\lhead{INF-324 -- Unidad 2}\rhead{Administración y PL/SQL}\cfoot{\thepage}
\setlength{\headheight}{14pt}

\lstdefinelanguage{OracleSQL}{
  morekeywords={CREATE,OR,REPLACE,PROCEDURE,FUNCTION,TRIGGER,SEQUENCE,RETURN,IS,AS,
  DECLARE,BEGIN,END,EXCEPTION,WHEN,THEN,ELSE,ELSIF,IF,FOR,IN,LOOP,WHILE,EXIT,
  LOCK,TABLE,ROW,SHARE,EXCLUSIVE,MODE,INSERT,INTO,VALUES,UPDATE,SET,DELETE,FROM,
  WHERE,SELECT,COMMIT,ROLLBACK,SAVEPOINT,RAISE,RAISE_APPLICATION_ERROR,DROP,
  BEFORE,AFTER,ON,EACH,NEW,OLD,START,WITH,INCREMENT,BY,MINVALUE,MAXVALUE,CYCLE,
  NOCYCLE,CACHE,NEXTVAL,CURRVAL,OPEN,FETCH,CLOSE,NO_DATA_FOUND,OTHERS},
  sensitive=false,morecomment=[l]{--},morestring=[b]'}
\lstset{language=OracleSQL,backgroundcolor=\color{grisCodigo},basicstyle=\ttfamily\small,
  keywordstyle=\color{sqlKeyword}\bfseries,stringstyle=\color{sqlString},
  commentstyle=\color{sqlComment}\itshape,numbers=left,numberstyle=\tiny\color{gray},
  numbersep=8pt,frame=single,rulecolor=\color{gray!35},breaklines=true,
  showstringspaces=false,upquote=true,tabsize=2,columns=fullflexible}

\newtcolorbox{objetivo}{colback=azulClaro,colframe=azulUCM,title=Objetivo de lectura,fonttitle=\bfseries}
\newtcolorbox{pausa}{colback=amarilloClaro,colframe=orange!70!black,title=Pausa activa,fonttitle=\bfseries}
\newtcolorbox{clave}{colback=verdeClaro,colframe=green!45!black,title=Idea clave,fonttitle=\bfseries}
\newtcolorbox{errorbox}{colback=rojoClaro,colframe=red!60!black,title=Error frecuente,fonttitle=\bfseries}
\newtcolorbox{rutina}{colback=blue!3,colframe=azulUCM,title=Rutina diaria de estudio,fonttitle=\bfseries}
\newcommand{\campo}[1]{\texttt{#1}}

\title{\textbf{Guía de lectura activa de la Unidad 2}\\[0.3em]
\large Administración de bases de datos y programación PL/SQL\\[0.5em]
\normalsize Preparación progresiva para INF-324}
\author{Construida a partir de las seis presentaciones y cinco guiones SQL de U2}
\date{Octubre de 2026}

\begin{document}
\maketitle\thispagestyle{empty}
\begin{center}\NombreEstudiante \hfill \textbf{Fecha de prueba:} \FechaPrueba\end{center}

\begin{objetivo}
Al terminar esta guía deberías poder explicar cómo se protege una base de datos, escribir bloques
PL/SQL reutilizables, controlar transacciones y concurrencia, y elegir entre una función, una
secuencia, un cursor o un trigger según el problema. La meta es comprender cada decisión antes de
copiar una solución.
\end{objetivo}

\vfill
\begin{center}\fbox{\parbox{0.9\textwidth}{
\textbf{Cómo usarla.} Lee una sección breve, tapa el código, reconstruye la idea con tus palabras y
responde las preguntas antes de mirar la pauta. En cada módulo marca los errores de concepto,
sintaxis y transacción por separado.\
}}
\end{center}
\vfill
\begin{flushleft}
\textbf{Fuentes locales revisadas:}
\begin{itemize}
\item Presentaciones en \campo{U2/presentaciones}: \campo{Clase 1}, \campo{Clase 1-2}, \campo{Clase 2}, \campo{Clase 2-2}, \campo{Clase 3} y \campo{Clase 4}.
\item Guiones en \campo{U2/codigos}: \campo{U2\_C1.SQL}, \campo{U2\_C1-2.SQL}, \campo{U2\_C2.SQL}, \campo{U2\_C2-2.SQL} y \campo{U2\_C4.SQL}.
\end{itemize}
\end{flushleft}
\newpage\tableofcontents\newpage

\section{Mapa de la unidad}
La Unidad 2 desarrolla el resultado de aprendizaje de administración de bases de datos mediante
una secuencia práctica. Las parejas de presentación y código se leen como una sola lección: primero
identifica el concepto, luego sigue el orden del script y finalmente modifica el ejemplo.

\begin{longtable}{p{0.13\textwidth} p{0.28\textwidth} p{0.45\textwidth}}
\toprule\textbf{Bloque} & \textbf{Fuentes} & \textbf{Pregunta guía}\\\midrule\endhead
1 & Clase 1 + \campo{U2\_C1.SQL} & ¿Cómo se construye y ejecuta un procedimiento almacenado?\\
2 & Clase 1-2 + \campo{U2\_C1-2.SQL} & ¿Cómo se coordinan usuarios, bloqueos, COMMIT, ROLLBACK y excepciones?\\
3 & Clase 2-2 + \campo{U2\_C2-2.SQL} & ¿Cuándo conviene una función y qué restricciones tiene?\\
4 & Clase 2 + \campo{U2\_C2.SQL} & ¿Cómo una secuencia genera claves sin depender de una tabla?\\
5 & Clase 3 & ¿Cómo se recorre el resultado de un SELECT con cursores?\\
6 & Clase 4 + \campo{U2\_C4.SQL} & ¿Qué evento debe activar un trigger y qué dato debe completar?\\
\bottomrule
\end{longtable}

\begin{rutina}
Cada sesión dura aproximadamente \MinutosSesion{} minutos: 5 para recordar sin mirar, 10 para leer,
20 para escribir o modificar SQL y 5 para corregir. Repite al día siguiente la pregunta en la que
obtuviste una respuesta roja. Una semana de práctica puede seguir el orden de los seis bloques y
cerrar con una simulación integral.
\end{rutina}

\section{Bloque 1: seguridad, PL/SQL y procedimientos}
\subsection{Lee y explica}
La seguridad de la información se trabaja como confidencialidad, disponibilidad e integridad. En
Oracle, PL/SQL agrega variables, estructuras de control y excepciones a SQL. Un bloque tiene las
secciones DECLARE, BEGIN, EXCEPTION y END; un procedimiento es un subprograma nombrado, compilado y
almacenado que puede recibir parámetros y ejecutar DML.

\begin{lstlisting}
CREATE OR REPLACE PROCEDURE CACB_CARGA_USUARIOS(
  COD_USUARIO_P NUMBER, NOMBRE_USUARIO_P VARCHAR2)
IS
BEGIN
  INSERT INTO CACB_USUARIO(COD_USUARIO, NOMBRE_USUARIO)
  VALUES(COD_USUARIO_P, NOMBRE_USUARIO_P);
END;
\end{lstlisting}

\begin{clave}Un procedimiento realiza una acción y puede recibir parámetros; una función debe
devolver un valor. No confundas el bloque anónimo que lo invoca con el procedimiento almacenado.\end{clave}
\begin{pausa}
Sin mirar el ejemplo: explica qué parte declara parámetros, qué parte ejecuta el INSERT y cómo
invocarías el procedimiento desde un bloque BEGIN...END.
\end{pausa}

\subsection{Preguntas}
\begin{enumerate}[leftmargin=*]
\item (Fácil) Nombra los tres objetivos de la seguridad de la información.
\item (Fácil) ¿Qué diferencia hay entre DECLARE y BEGIN?
\item (Fácil) ¿Qué tres operaciones DML se trabajan en los procedimientos de U2?
\item (Intermedia) ¿Por qué CREATE OR REPLACE es útil al practicar?
\item (Intermedia) Escribe un procedimiento que elimine un usuario por código.
\item (Difícil) Diseña un procedimiento I/U/D que reciba una opción y justifica qué parámetros pueden tener valor NULL.
\end{enumerate}
\ifmostrarSoluciones\begin{clave}\textbf{Pauta breve:} confidencialidad, disponibilidad e integridad. DECLARE prepara variables y BEGIN ejecuta. DML: INSERT, UPDATE y DELETE. CREATE OR REPLACE recompila y reemplaza la versión anterior. El procedimiento I/U/D debe validar la opción antes de ejecutar la operación.\end{clave}\fi

\section{Bloque 2: concurrencia, transacciones y excepciones}
\subsection{Lee y explica}
La concurrencia aparece cuando varios usuarios comparten recursos. \campo{LOCK TABLE} impone un
modo de bloqueo; \campo{ROW EXCLUSIVE} permite trabajar por filas y evita ciertos bloqueos
incompatibles. \campo{COMMIT} confirma, \campo{ROLLBACK} deshace y \campo{SAVEPOINT} ofrece un
punto intermedio. Las excepciones se controlan después de BEGIN y pueden producirse con \campo{RAISE}
o \campo{RAISE\_APPLICATION\_ERROR}.

\begin{lstlisting}
BEGIN
  LOCK TABLE CACB_USUARIO IN ROW EXCLUSIVE MODE;
  INSERT INTO CACB_USUARIO VALUES (1, 'CARLOS');
  COMMIT;
EXCEPTION
  WHEN OTHERS THEN
    RAISE_APPLICATION_ERROR(-20010, 'ERROR DE CARGA');
    ROLLBACK;
END;
\end{lstlisting}

\begin{errorbox}Un SELECT no bloquea una tabla. Además, un bloqueo permanece hasta COMMIT, ROLLBACK o un retroceso a un SAVEPOINT válido. Revisa el flujo de excepción y la posición de cada instrucción.\end{errorbox}
\subsection{Preguntas}
\begin{enumerate}[leftmargin=*]
\item (Fácil) ¿Qué instrucción confirma una transacción?
\item (Fácil) ¿Qué instrucción deshace el trabajo actual?
\item (Fácil) ¿Qué propósito tiene LOCK TABLE?
\item (Intermedia) Compara ROW SHARE, ROW EXCLUSIVE y EXCLUSIVE.
\item (Intermedia) ¿Por qué una tabla bloqueada aún puede ser consultada?
\item (Difícil) Ordena un procedimiento de carga que bloquee, inserte, confirme y revierta al fallar.
\end{enumerate}
\ifmostrarSoluciones\begin{clave}\textbf{Pauta breve:} COMMIT confirma y ROLLBACK deshace. El modo de bloqueo determina qué operaciones simultáneas se permiten. Las consultas no colocan bloqueo de tabla. Un flujo robusto pone LOCK antes del DML, COMMIT al terminar y ROLLBACK en el manejador de excepciones.\end{clave}\fi

\section{Bloque 3: funciones PL/SQL}
Una función es un bloque nombrado que devuelve un valor y puede usarse dentro de una expresión o
consulta. El código de U2 usa \campo{SELECT INTO} para llevar una columna a una variable y captura
\campo{NO\_DATA\_FOUND}. Las funciones llamadas desde SQL no deben ejecutar INSERT, UPDATE o DELETE.

\begin{lstlisting}
CREATE OR REPLACE FUNCTION CACB_OBTENER_SUELDO(EMPLOYEE_ID_F NUMBER)
RETURN NUMBER IS
  SUELDO_OBTENIDO NUMBER := 0;
BEGIN
  SELECT SALARY INTO SUELDO_OBTENIDO
  FROM OEHR_EMPLOYEES WHERE EMPLOYEE_ID = EMPLOYEE_ID_F;
  RETURN SUELDO_OBTENIDO;
EXCEPTION
  WHEN NO_DATA_FOUND THEN RETURN -1;
END;
\end{lstlisting}

\begin{pausa}Tapa el código y explica qué ocurre si el sueldo es NULL y qué ocurre si no existe el empleado. Luego indica cómo llamarías la función desde DUAL.\end{pausa}
\subsection{Preguntas}
\begin{enumerate}[leftmargin=*]
\item (Fácil) ¿Qué debe declarar RETURN?
\item (Fácil) ¿Para qué sirve SELECT INTO?
\item (Fácil) ¿Qué excepción indica que no se encontró una fila?
\item (Intermedia) Diferencia una función SQL integrada de una función escrita en PL/SQL.
\item (Intermedia) ¿Por qué una función de consulta no debe ejecutar DML?
\item (Difícil) Diseña una función que devuelva 0 si un valor es NULL y -1 si no existe la fila.
\end{enumerate}
\ifmostrarSoluciones\begin{clave}\textbf{Pauta breve:} RETURN declara el tipo y devuelve el valor. SELECT INTO asigna el resultado a una variable. NO_DATA_FOUND indica ausencia de fila. La función de ejemplo usa IF para transformar NULL en 0 y un manejador para devolver -1.\end{clave}\fi

\section{Bloque 4: secuencias}
Una secuencia genera enteros únicos para uno o varios usuarios y puede alimentar una clave primaria.
Sus números avanzan independientemente de COMMIT o ROLLBACK, por lo que pueden quedar saltos.
\campo{NEXTVAL} incrementa y devuelve el siguiente valor; \campo{CURRVAL} devuelve el valor actual
después de inicializar la secuencia.

\begin{lstlisting}
CREATE SEQUENCE CACB_PK_USUARIO
  START WITH 1 INCREMENT BY 1 MINVALUE 1
  MAXVALUE 2500 CACHE 200 NOCYCLE;

INSERT INTO CACB_USUARIO(COD_USUARIO, NOMBRE_USUARIO)
VALUES(CACB_PK_USUARIO.NEXTVAL, 'CARLOS');
\end{lstlisting}

\begin{clave}START WITH fija el inicio; INCREMENT BY fija el paso; MINVALUE y MAXVALUE limitan; CYCLE permite reiniciar y NOCYCLE evita reutilizar números.\end{clave}
\subsection{Preguntas}
\begin{enumerate}[leftmargin=*]
\item (Fácil) ¿Qué pseudocolumna se usa primero?
\item (Fácil) ¿Qué diferencia hay entre NEXTVAL y CURRVAL?
\item (Fácil) ¿Por qué una secuencia puede presentar huecos?
\item (Intermedia) Interpreta START WITH 1, INCREMENT BY 1 y MAXVALUE 2500.
\item (Intermedia) ¿Qué efecto tiene NOCYCLE al alcanzar MAXVALUE?
\item (Difícil) Explica por qué una secuencia es preferible a MAX(id)+1 en un procedimiento concurrente.
\end{enumerate}
\ifmostrarSoluciones\begin{clave}\textbf{Pauta breve:} NEXTVAL inicializa e incrementa; CURRVAL lee el actual. Los huecos aparecen por consumos concurrentes o transacciones revertidas. NOCYCLE detiene la generación al límite. La secuencia evita colisiones entre sesiones.\end{clave}\fi

\section{Bloque 5: cursores PL/SQL}
Los cursores permiten manejar el conjunto de filas devuelto por SELECT. Los implícitos los gestiona
Oracle para resultados de una fila; los explícitos se declaran y recorren cuando se necesita procesar
varias filas o reutilizar la consulta. La declaración del cursor no incluye INTO; el orden se indica
con ORDER BY cuando es necesario.

\begin{lstlisting}
DECLARE
  CURSOR c_usuarios IS
    SELECT COD_USUARIO, NOMBRE_USUARIO
    FROM CACB_USUARIO ORDER BY COD_USUARIO;
BEGIN
  FOR r IN c_usuarios LOOP
    DBMS_OUTPUT.PUT_LINE(r.NOMBRE_USUARIO);
  END LOOP;
END;
\end{lstlisting}

\begin{pausa}Dibuja mentalmente el conjunto de filas, el registro r y el momento en que el FOR abre, obtiene y cierra el cursor.\end{pausa}
\subsection{Preguntas}
\begin{enumerate}[leftmargin=*]
\item (Fácil) ¿Qué tipo de cursor se usa para varias filas?
\item (Fácil) ¿Qué representa un cursor internamente?
\item (Fácil) ¿Dónde debe ir ORDER BY si importa el orden?
\item (Intermedia) ¿Por qué no se incluye INTO en la declaración?
\item (Intermedia) Compara FOR r IN cursor con OPEN, FETCH y CLOSE manuales.
\item (Difícil) Propón un cursor que recorra usuarios y actualice una tabla de auditoría, indicando dónde controlarías errores.
\end{enumerate}
\ifmostrarSoluciones\begin{clave}\textbf{Pauta breve:} el cursor explícito maneja varias filas. Es un segmento de memoria asociado al resultado. INTO pertenece a la obtención de una fila, no a la declaración. FOR gestiona automáticamente apertura, FETCH y cierre.\end{clave}\fi

\section{Bloque 6: triggers e integridad}
Un trigger es un bloque PL/SQL asociado a una tabla que se ejecuta por un evento DML. Sus componentes
son nombre, evento, momento y cuerpo; \campo{FOR EACH ROW} lo hace ejecutar por fila. \campo{:NEW}
y \campo{:OLD} permiten leer o modificar valores de la fila. Un trigger no puede ejecutar COMMIT,
ROLLBACK ni SAVEPOINT.

\begin{lstlisting}
CREATE OR REPLACE TRIGGER CACB_PROMEDIO_AUTO
BEFORE INSERT ON CACB_DETALLE_NOTAS
FOR EACH ROW
BEGIN
  :NEW.PROMEDIO := ROUND((:NEW.NOTA1 + :NEW.NOTA2 + :NEW.NOTA3)/3, 1);
END;
\end{lstlisting}

\begin{errorbox}No coloques COMMIT dentro del trigger. El ejemplo calcula un dato derivado antes del INSERT; la transacción la controla el bloque o procedimiento que ejecuta la operación.\end{errorbox}
\subsection{Preguntas}
\begin{enumerate}[leftmargin=*]
\item (Fácil) ¿Qué operación activa el trigger del ejemplo?
\item (Fácil) ¿Qué representa :NEW?
\item (Fácil) ¿Qué cláusula hace que se ejecute por fila?
\item (Intermedia) Diferencia BEFORE INSERT de AFTER INSERT.
\item (Intermedia) Menciona dos usos de triggers además de calcular promedios.
\item (Difícil) Diseña un trigger que normalice a mayúsculas un nombre antes de insertarlo y explica cómo tratarías NULL.
\end{enumerate}
\ifmostrarSoluciones\begin{clave}\textbf{Pauta breve:} el evento es INSERT, :NEW contiene los valores nuevos y FOR EACH ROW ejecuta una vez por fila. Los triggers también mantienen integridad y apoyan auditoría. La normalización puede asignar UPPER(:NEW.nombre) si el valor no es NULL.\end{clave}\fi

\section{Simulación integradora}
Resuelve sin mirar los módulos y luego revisa cada decisión.
\begin{enumerate}[leftmargin=*]
\item Crea una tabla de usuarios con clave primaria.
\item Crea una secuencia para la clave y justifica MAXVALUE, CACHE y NOCYCLE.
\item Escribe un procedimiento de carga que use NEXTVAL, LOCK TABLE, COMMIT y ROLLBACK.
\item Agrega una función que devuelva un dato derivado y trate NO_DATA_FOUND.
\item Declara un cursor explícito para recorrer los usuarios ordenados.
\item Agrega un trigger que complete un campo antes de insertar, sin controlar la transacción.
\end{enumerate}

\begin{longtable}{p{0.36\textwidth} *{3}{>{\centering\arraybackslash}p{0.16\textwidth}}}
\toprule\textbf{Habilidad} & \textbf{Intento 1} & \textbf{Intento 2} & \textbf{Intento 3}\\\midrule
Procedimiento y bloque PL/SQL & & & \\
Bloqueo y transacción & & & \\
Excepciones & & & \\
Funciones & & & \\
Secuencias & & & \\
Cursores & & & \\
Triggers & & & \\
\bottomrule
\end{longtable}

\section*{Checklist final}
\begin{itemize}[leftmargin=*]
\item Puedo explicar confidencialidad, disponibilidad e integridad.
\item Puedo distinguir procedimiento, función, secuencia, cursor y trigger.
\item Puedo ordenar LOCK, DML, COMMIT y ROLLBACK en una transacción.
\item Puedo capturar NO_DATA_FOUND y generar un error propio.
\item Puedo explicar NEXTVAL, CURRVAL y los límites de una secuencia.
\item Puedo recorrer filas con un cursor y evitar INTO en su declaración.
\item Puedo usar :NEW/:OLD y mantener el control de transacciones fuera del trigger.
\end{itemize}

\end{document}
